\documentclass{article}

% Uncomment for final submission:
% \usepackage[final]{aaj2026}

\bibliographystyle{plainnat}
\usepackage[numbers]{natbib}
\usepackage[utf8]{inputenc}
\usepackage[T1]{fontenc}
\usepackage[hidelinks]{hyperref}
\usepackage{url}
\urlstyle{same}
\usepackage{booktabs}
\usepackage{amsfonts}
\usepackage{nicefrac}
\usepackage{microtype}
\usepackage{xcolor}
\usepackage{amsmath}
\usepackage{mathtools}
\usepackage{graphicx}
\graphicspath{{../plots/figures/}}

\title{Fast Monte Carlo Pricing of Asian Options under Rough Fractional Stochastic Volatility}

\author{%
  Tiago Flora \\
  Department of Mathematics \\
  University of Chicago \\
  \texttt{tflora@uchicago.edu}
}

\begin{document}

\maketitle

\begin{abstract}
We study Monte Carlo pricing of a fixed-strike arithmetic Asian call option under the
Rough Fractional Stochastic Volatility (RFSV) model~\citep{gatheral2018volatility},
in which log-volatility follows a fractional Brownian motion with Hurst exponent
$H = 0.10$.  The non-Markovian structure of fBM forces every sample path to be drawn
from a dense $N \times N$ Gaussian distribution, requiring an $O(N^3)$ Cholesky setup
and an $O(MN^2)$ per-path simulation cost.  We implement and benchmark an exact dense
Cholesky baseline against two algorithms that exploit structure in the covariance
matrix: a circulant-embedding FFT method ($O(N \log N)$ per path, exploiting the
stationarity of fBM increments), and a global randomized SVD method ($O(Nk)$ per
path, exploiting low-rank structure).  We show that the circulant embedding is exact
for every $H \leq \nicefrac{1}{2}$, with smallest eigenvalue
$\Delta t^{2H}(N^{2H} - (N-1)^{2H})$ in closed form.  Over
$N \in \{252, 500, 1000\}$ the Cholesky and FFT costs scale as $N^{1.33}$ and
$N^{0.95}$, but at $N = 1000$ the FFT method is only $\approx 9\%$ faster than a
triangular Cholesky baseline, because random number generation rather than the
covariance factor dominates its per-path cost.  The rank-32 rSVD method is
$2.7\times$ faster than Cholesky at $N = 1000$, but its covariance error decays
slowly, from $1.9\%$ at $k = 32$ to $1.2\%$ at $k = 128$, because the rough kernel's
singularity lies on the diagonal, where a global low-rank approximation cannot
isolate it.
\end{abstract}

\section{Introduction}

\subsection{The option contract}

I am interested in pricing fixed-strike, arithmetic-average Asian call options.  The
payoff at expiry $T$ is
\begin{equation}
  C(T) = \max\!\left(A(0,T) - K,\, 0\right), \qquad
  A(0,T) = \frac{1}{N}\sum_{j=1}^{N} S_{t_j},
\end{equation}
where $S_{t_j}$ is the price of the underlying at time $t_j = jT/N$ and $K$ is the
strike.  Path-dependence through the average $A$ prevents the Black--Scholes machinery
from yielding an exact closed form, making approximations or simulation necessary.
\citet{levy1992pricing} proposed a moment-matching approximation: matching the first
two moments of the discrete arithmetic average to a lognormal yields a tractable
formula that is accurate under constant-volatility Black--Scholes dynamics, where $S_t$
is lognormal.

The situation is more difficult under Rough Fractional Stochastic Volatility (RFSV)
\citep{gatheral2018volatility}, in which $\log\sigma_t = \nu W_t^H$, where $W_t^H$ is
a fractional Brownian motion with Hurst exponent $H \approx 0.10$.  Under RFSV, $S_t$
is no longer lognormal, and the two-moment Lévy approximation cannot capture the heavy
skew and fat tails introduced by a large vol-of-vol $\nu$.  Computing even the second
joint moment $\mathbb{E}[S_{t_i}S_{t_j}]$ analytically is intractable: it requires the
expectation of the integrated volatility path, which depends on the full continuous
history of $W^H$ up to each time point.

Monte Carlo sidesteps the distributional intractability by estimating the expectation
directly from simulated paths.  By the law of large numbers,
\begin{equation}
  \frac{1}{M}\sum_{m=1}^{M} e^{-rT}\max\!\left(A^{(m)} - K,\,0\right)
  \xrightarrow{M \to \infty}
  \mathbb{E}^{\mathbb{Q}}\!\left[e^{-rT}\max(A - K,\, 0)\right],
\end{equation}
where $\mathbb{Q}$ is the risk-neutral measure and $A^{(m)}$ is the arithmetic average
along the $m$-th simulated path.  No closed-form density or moment-matching is required;
we need only the ability to simulate correlated log-volatility paths.

\subsection{Why rough volatility}

\citet{gatheral2018volatility} provide compelling empirical evidence that equity
log-realized-variance behaves as fractional Brownian motion with $H \approx 0.10$.
Estimating $H$ by ordinary least squares on log-realized-variance increments across
five major indices (S\&P 500, NASDAQ, DAX, FTSE 100, Nikkei 225) and time scales
ranging from one day to one month, they find a consistent Hurst exponent well below
$\nicefrac{1}{2}$, with $R^2 > 0.90$ in each case.  I take $H = 0.10$ and their
vol-of-vol estimate from high-frequency realized variance data \citep{oxfordman},
$\nu \approx 0.3$.  One subtlety concerns units: \citet{gatheral2018volatility}
measure time in days, while my engine measures it in years ($T = 1$).  Because the
log-vol increment variance scales as $\nu^2\Delta^{2H}$, the same fit in year units is
$\nu = 0.3 \cdot 252^{0.1} \approx 0.52$, which is the value I use.  My own variogram fit to
non-overlapping 5-day realized variance from \texttt{yfinance}, with time in years,
gives $H \approx 0.12$ ($R^2 \approx 0.98$ across a handful of lags) and
$\nu \approx 0.87$, the latter inflated by measurement noise in weekly
squared-return variance.

Standard Brownian motion corresponds to $H = \nicefrac{1}{2}$, with uncorrelated increments. 
At $H = 0.10$, volatility paths are far rougher than BM: increments
are negatively correlated at all lags, the path is less regular than classical
stochastic volatility models, and the realized variance process mean-reverts on
short timescales. The short-dated implied volatility smile
steepens beyond what classical stochastic volatility models expect, and
volatility clustering manifests on timescales an order of magnitude shorter than, e.g., Heston.

The RFSV model I studied is \citep{gatheral2018volatility}
\begin{equation}
  \log\sigma_t = \log\sigma_0 + \nu W_t^H, \qquad
  dS_t = r S_t\,dt + \sigma_t S_t\, dW_t^{\mathbb{Q}},
\end{equation}
where $W_t^H$ is the fractional Brownian motion of \citet{mandelbrot1968fractional}
with covariance $\mathbb{E}[W_t^H W_s^H] = \tfrac{1}{2}(t^{2H}+s^{2H}-|t-s|^{2H})$.
I set $r = 0$ throughout for simplicity (consistent with \texttt{params.hpp}); the
drift term drops out of all pricing formulas but is retained in the SDE for
completeness.
I use $H = 0.10$, $\nu = 0.52$, and base volatility $\sigma_0 = 0.2$ throughout, and
also explore the sensitivity of option prices to $H$ and $\nu$ in
Section~\ref{sec:sensitivity}.  The base level is not a cosmetic choice.  An earlier
version of this study used $\sigma_0 = 1$ (100\% volatility); at $\nu = 0.52$ that
payoff is so heavy-tailed that a single path in $10^6$ carried 34\% of the sample
variance, and Monte Carlo error bars were meaningless.  At $\sigma_0 = 0.2$ the largest
path carries $0.11\%$, and the payoff standard deviation is stable across batches.
Strictly, the payoff variance is infinite for any lognormal volatility model, because
$\mathbb{E}[S_t^2]$ contains $\mathbb{E}[\exp(\Delta t\, e^{2\nu W})]$; in practice the
divergence sits so far in the tail that it is never sampled.

\subsection{The computational bottleneck}
\label{sec:bottleneck}

Naïve Monte Carlo for RFSV requires $O(N^3)$ setup work.  Because fBM is
non-Markovian, there is no forward recursion: the value of $W_t^H$ depends on the
entire history of the driving noise, not just its most recent value.  Generating each
sample path therefore requires drawing from the full $N$-dimensional Gaussian
$\mathcal{N}(0, C)$, where $C$ is a dense $N \times N$ positive-definite matrix with
$C_{ij} = \tfrac{1}{2}(t_i^{2H}+t_j^{2H}-|t_i-t_j|^{2H})$.  The standard approach
factorizes $C = LL^\top$ via Cholesky decomposition once ($O(N^3/3)$ flops), then
generates each path by computing $\nu Lz$ for $z \sim \mathcal{N}(0,I)$ ($O(N^2)$
flops per path, a lower-triangular matrix-vector multiply).

At $N = 252$ trading days and $M = 10{,}000$ paths the costs are concrete.  Cholesky
decomposition requires exactly $N^3/3 \approx 5.3$ million floating-point operations
(the leading constant $\nicefrac{1}{3}$ compares favourably to $\nicefrac{2}{3}$ for
LU); this is a one-time cost.  The per-path triangular multiply costs $O(N^2) \approx 63{,}000$ flops, so
the total Monte Carlo loop costs $O(MN^2) \approx 630$ million flops, more than
$100\times$ the factorization.  MC pricing error falls as $O(M^{-1/2})$: with
$\sigma_\mathrm{payoff} \approx 9.5$ (Section~\ref{sec:convergence}), $M = 10{,}000$
paths give a standard error of $\approx 0.095$ price units, about $1.8\%$ of the
at-the-money price of $\approx 5.3$, and halving it requires four times as many paths.

Trading desks building sensitivity surfaces run parameter sweeps over $H$, $\nu$, and
$K$, easily multiplying $M$ by factors of $10$--$100$.  Intraday risk monitoring at
$N \approx 19{,}656$ (five-minute intervals over 252 trading days of 6.5 hours each:
$252 \times 78 = 19{,}656$) increases the per-path cost by $\sim 6{,}000\times$
relative to daily steps.

\subsection{Approach and scope}
\label{sec:scope}

I implement and benchmark three algorithms that reduce the cost of fBM path simulation
by exploiting structure in the covariance matrix $C$:

\begin{enumerate}
  \item \textbf{Dense Cholesky} ($O(N^2)$/path, exact): exploits the
        positive-definiteness of $C$ to factor it exactly.  By ``exact simulation''
        I mean that the method is free of time-discretization bias: it samples the correct
        fBM distribution for any finite $N$, unlike Euler--Maruyama approximations.
        Serves as the correctness baseline against which the other methods are validated.

  \item \textbf{Circulant embedding \& FFT} ($O(N\log N)$/path): exploits the
        fact that the increments of fBM (fractional Gaussian noise, fGn) are
        stationary. Their covariance matrix is Toeplitz, which can be embedded
        into a circulant matrix, which can be diagonalized by the DFT
        \citep{davies1987tests, wood1994simulation}.  For $H \leq \nicefrac{1}{2}$
        the embedding is guaranteed positive \citep{craigmile2003simulating}, so the
        method is exact on the grid.

  \item \textbf{Global low-rank rSVD} ($O(Nk)$/path): exploits the low-rank
        structure of smooth off-diagonal blocks of $C$: the off-diagonal blocks admit low-rank
        approximations.  A randomized SVD \citep{halko2011finding} compresses $C
        \approx U\,\mathrm{diag}(S)\,U^\top$ at rank $k$, reducing per-path cost from
        $O(N^2)$ to $O(Nk)$.  In practice the implementation is a global
        approximation of the entire $N\times N$ matrix (including the singular
        near-diagonal region) rather than of the smooth off-diagonal blocks alone.
        This is blunter, and must allocate rank budget to the diagonal singularity,
        which is why the sampler still discards $13\%$ of the path variance at
        $k = 128$ for $H = 0.10$, and underprices by $\approx 4\%$
        (§\ref{sec:errorrank}).
\end{enumerate}

Section~\ref{sec:foundations} develops the mathematical properties of $C$ that each
method exploits.  Section~\ref{sec:algorithms} describes each algorithm and its
complexity.  Sections~\ref{sec:benchmarks} and~\ref{sec:validation} present empirical
timing, memory, accuracy, and model validation results.

The project does not address variance reduction (quasi-Monte Carlo, control variates),
model extensions beyond fBM-RFSV (rough Bergomi, multifactor), or GPU parallelism,
though Section~\ref{sec:futurework} situates each of these as a natural next step.

\section{Mathematical Foundations}
\label{sec:foundations}

\subsection{Fractional Brownian motion and its covariance structure}

Fractional Brownian motion $W_t^H$ is the unique (up to scaling) continuous Gaussian
process with stationary increments and the self-similarity property
$W_{ct}^H \stackrel{d}{=} c^H W_t^H$ for all $c > 0$
\citep{mandelbrot1968fractional}.  Its covariance kernel is
\begin{equation}
  \label{eq:fbm-cov}
  \mathbb{E}[W_t^H W_s^H]
  = \tfrac{1}{2}\!\left(t^{2H} + s^{2H} - |t-s|^{2H}\right),
\end{equation}
which is symmetric, positive-definite for $H \in (0,1)$, and reduces to the standard
Brownian covariance $\min(t,s)$ at $H = \nicefrac{1}{2}$.  The self-similarity
$W_{ct}^H \stackrel{d}{=} c^H W_t^H$ and the stationarity of increments
$W_{t+h}^H - W_t^H \stackrel{d}{=} W_h^H$ are the properties that make fBM a
tractable model of rough volatility; §\ref{sec:toeplitz} exploits increment
stationarity directly.  The critical property for simulation is that fBM is
non-Markovian for $H \neq \nicefrac{1}{2}$: the value of $W_t^H$ depends on
the entire history of the driving noise, not just its most recent value.  This rules
out any forward recursion and forces each sample path to be drawn from a joint
$N$-dimensional Gaussian.

Discretizing $[0,T]$ into $N$ equally spaced points $t_j = jT/N$ yields an $N\times N$
covariance matrix $C$ with $C_{jk} = \mathbb{E}[W_{t_j}^H W_{t_k}^H]$ as in
equation~\eqref{eq:fbm-cov}.  The matrix $C$ is dense (every pair of time points is correlated) and positive-definite
for all $H \in (0,1)$ and $N \geq 1$.  This density explains the $O(N^2)$ per-path
cost: sampling from
$\mathcal{N}(0,\nu^2 C)$ by any exact method requires touching all $N^2/2$ entries
of $C$.  Algorithm~1 (§\ref{sec:cholesky}) accepts this cost and factorizes $C$
exactly; Algorithms~2 and~3 avoid it by exploiting structure that $C$ possesses.

\subsection{Stationarity of fGn and the Toeplitz property}
\label{sec:toeplitz}

Although the fBM covariance matrix is dense, its increment process has a
crucial structural simplification.  Define the fractional Gaussian noise (fGn)
increments $\delta W_m = W_{(m+1)\Delta t}^H - W_{m\Delta t}^H$ for $m = 0, \ldots,
N-1$, where $\Delta t = T/N$.  Stationarity of fBM increments implies
\begin{equation}
  \label{eq:fgn-cov}
  \gamma(k) \coloneqq \mathrm{Cov}(\delta W_m,\, \delta W_{m+k})
  = \frac{\Delta t^{2H}}{2}\bigl(|k+1|^{2H} + |k-1|^{2H} - 2|k|^{2H}\bigr),
\end{equation}
which depends only on the lag $k = |m - n|$, not on the absolute index $m$
\citep{wood1994simulation}.  The $N\times N$ covariance matrix of
$(\delta W_0, \ldots, \delta W_{N-1})$ is therefore Toeplitz: its $(m,n)$
entry equals $\gamma(|m-n|)$, so every diagonal is constant.

Figure~\ref{fig:structure} illustrates the contrast.  Panels~(a) and~(b) show the
$N\times N$ covariance heatmaps of fBM and fGn respectively.  The fBM matrix is not
Toeplitz: row $m$ of $C$ varies with $m$ because the variance $t_m^{2H}$ grows with
time.  The fGn matrix has constant diagonals, confirming the Toeplitz property exactly.
Panel~(c) shows diagonal slices $C[m, m+k]$ vs.\ lag $k$ for two different rows
$m$: the fGn curves overlap to numerical precision, while the fBM curves diverge.

\begin{figure}[t]
  \centering
  \includegraphics[width=\linewidth]{structure_analysis.png}
  \caption{Structural properties of the fBM and fGn covariance matrices ($H=0.10$).
    \textbf{(a)} fBM covariance $C(t_m, t_n)$: dense, not Toeplitz (variance grows
    with time).
    \textbf{(b)} fGn covariance $\gamma(|m-n|)$: Toeplitz (constant diagonals).
    \textbf{(c)} Diagonal slice $C[m, m+k]$ vs.\ lag $k$: fGn rows for $m=0$ and
    $m=N/4$ overlap (Toeplitz); fBM rows diverge (non-stationary).
    \textbf{(d)} Normalised singular values $\sigma_k/\sigma_1$ of the off-diagonal
    block (solid) and normalised eigenvalues of the full matrix $C$ (dashed) at
    $H=0.10$ and $H=0.50$.  The off-diagonal block compresses to rank $\leq 3$ at
    both $H$; the full spectrum decays slowly, most slowly at $H=0.10$.
    Controlled: $H=0.10$, $N=64$ (heatmaps), $N=128$ (SVD).
    Independent: process type or Hurst exponent.
    Dependent: covariance pattern and singular value magnitude.}
  \label{fig:structure}
\end{figure}

Algorithm~2 (§\ref{sec:fft}) depends on this Toeplitz structure: only Toeplitz
matrices embed into a circulant diagonalizable by the discrete Fourier transform.
Applying the FFT method directly to the fBM covariance matrix 
produces negative eigenvalues and pathological path output.

\subsection{Circulant embedding and its positivity condition}
\label{sec:circulant}

A Toeplitz matrix $T$ with first row $(\gamma(0), \gamma(1), \ldots, \gamma(N-1))$
can be embedded into a $2N\times 2N$ circulant matrix $C_\mathrm{circ}$ whose
first row is the symmetric extension
\begin{equation}
  c = \bigl(\gamma(0),\, \gamma(1),\, \ldots,\, \gamma(N-1),\, 0,\,
             \gamma(N-1),\, \ldots,\, \gamma(1)\bigr).
\end{equation}
Circulant matrices are diagonalized by the DFT: the eigenvalues of $C_\mathrm{circ}$
are exactly $\lambda = \mathrm{FFT}(c)$, computable in $O(N\log N)$ operations.
To sample, draw $W = a + ib$ with $a, b \sim \mathcal{N}(0, I_{2N})$ independent
(so each component has real and imaginary parts that are independent standard
normals),
scale component-wise by $\sqrt{\lambda_j / (2N)}$, apply the inverse FFT, and take
the real part of the first $N$ elements \citep{davies1987tests}.  One forward FFT (once per experiment) to compute
$\{\lambda_j\}$, then one inverse FFT per path: total $O(N\log N)$ per path.

The method is valid only if every eigenvalue $\lambda_j$ is non-negative, so that
$\sqrt{\lambda_j}$ is real \citep{wood1994simulation}.  For rough $H$ this holds at
every $N$, not just asymptotically.  \citet{craigmile2003simulating} proves that the
Davies--Harte embedding is non-negative definite for any stationary process whose
autocovariance is non-positive at all non-zero lags, a class that includes fGn with
$H < \nicefrac{1}{2}$.  The argument is short enough to give in full for the
embedding used here.  Because $x \mapsto x^{2H}$ is concave for $H < \nicefrac{1}{2}$,
the second difference in equation~\eqref{eq:fgn-cov} is negative, so $\gamma(k) < 0$
for every $k \geq 1$.  The eigenvalues of the symmetric circulant are
\begin{equation}
  \lambda_j = \gamma(0) + 2\sum_{k=1}^{N-1} \gamma(k)\cos\!\left(\frac{\pi jk}{N}\right),
  \qquad j = 0, \ldots, 2N-1,
\end{equation}
and because each $\gamma(k)$ is negative, replacing every cosine by $1$ can only lower
the sum.  The smallest eigenvalue is therefore $\lambda_0$, and the sum telescopes:
\begin{equation}
  \label{eq:lambda-min}
  \begin{split}
  \lambda_{\min} = \lambda_0
    &= \gamma(0) + 2\sum_{k=1}^{N-1}\gamma(k) \\
    &= \Delta t^{2H}\bigl(N^{2H} - (N-1)^{2H}\bigr)
     \approx 2H\,\Delta t^{2H} N^{2H-1} > 0.
  \end{split}
\end{equation}
The FFT method is thus exact on the grid for every $H \leq \nicefrac{1}{2}$ (at
$H = \nicefrac{1}{2}$ the increments are white noise and every $\lambda_j = \gamma(0)$).
The C++ pricer checks the sign of every eigenvalue and throws rather than clipping.

Equation~\eqref{eq:lambda-min} also shows that the margin is thin.  Relative to the
increment variance $\gamma(0) = \Delta t^{2H}$, the smallest eigenvalue is
$N^{2H} - (N-1)^{2H}$: $0.24\%$ at $H = 0.10$, $N = 252$ and $0.08\%$ at $N = 1000$,
decaying like $N^{-(1-2H)}$.  This mirrors the fGn spectral density, which vanishes
at zero frequency when $H < \nicefrac{1}{2}$.  Section~\ref{sec:stability} confirms
the formula numerically.  The key pitfall is $\gamma(0)$ itself: dropping the
$|k-1|^{2H}$ term at $k = 0$ halves it, which breaks the bound and spuriously makes
$28\%$ of the eigenvalues negative at $H = 0.10$.  An earlier version of my
stability analysis made exactly this mistake and concluded, wrongly, that the FFT
method needed eigenvalue clipping.

\subsection{Low-rank structure of smooth off-diagonal blocks}
\label{sec:lowrank}

The fBM kernel $C(t,s) = \tfrac{1}{2}(t^{2H}+s^{2H}-|t-s|^{2H})$ has a singularity at $t = s$, 
since the term $|t-s|^{2H}$ introduces one. For $H = 0.10$ this gives $|t-s|^{0.2}$, whose derivative
$0.2|t-s|^{-0.8}$ is unbounded as $t \to s$. The function is H\"older-continuous
with exponent $0.2$ but strictly non-differentiable on the diagonal.

Off-diagonal blocks of $C$ (submatrices coupling two disjoint time intervals) lie
in the smooth region $t \neq s$.  Smooth bivariate kernels generate matrices with
rapidly decaying singular values on such blocks \citep{candes2009fast}.

Figure~\ref{fig:structure}(d) confirms this, and locates the difficulty precisely.
The off-diagonal block coupling the first half of the grid to the second half is
exactly rank~1 at $H = 0.50$, where the kernel is $\min(t,s)$.  At $H = 0.10$ its
singular values still fall geometrically: below $1\%$ of $\sigma_1$ by rank $k = 3$
and below $10^{-4}$ by $k = 5$.  The spectrum of the full matrix $C$ behaves very
differently.  It decays only algebraically, and more slowly at $H = 0.10$ than at
$H = 0.50$.  The slow decay that limits a global low-rank approximation therefore
comes from the near-diagonal singularity, not from the far field.  This is exactly
the structure a hierarchical matrix exploits (§\ref{sec:hmatrix-future}).
\section{The three algorithms}
\label{sec:algorithms}

Table~\ref{tab:complexity} summarizes the three methods.  The rest of this section
describes each in turn.

\begin{table}[t]
  \caption{Complexity of the three fBM path generation methods.  $N$: time steps;
    $M$: paths; $k$: rSVD target rank.  ``Exact'' means the method samples the
    correct Gaussian distribution on the grid, up to floating-point error; for the
    FFT method this rests on the positivity result of §\ref{sec:circulant}.}
  \label{tab:complexity}
  \centering
  \begin{tabular}{lllll}
    \toprule
    Method & Setup & Per path & Memory & Accuracy \\
    \midrule
    Dense Cholesky   & $O(N^3)$        & $O(N^2)$        & $O(N^2)$  & Exact \\
    Circulant + FFT  & $O(N\log N)$    & $O(N\log N)$    & $O(N)$    & Exact \\
    Low-rank rSVD    & $O(N^2 k)$      & $O(Nk)$         & $O(Nk)$   & Approximate \\
    \bottomrule
  \end{tabular}
\end{table}

\subsection{Algorithm 1: Dense Cholesky}
\label{sec:cholesky}

The fBM covariance matrix $C$ is symmetric positive-definite for all $H \in (0,1)$
and $N \geq 1$, which is sufficient for an exact Cholesky factorization $C = LL^\top$.

\paragraph{Algorithm}
\begin{enumerate}
  \item Build $C \in \mathbb{R}^{N\times N}$ analytically from
        equation~\eqref{eq:fbm-cov}: $O(N^2)$ flops.
  \item Factorize $C = LL^\top$ in place via Cholesky decomposition
        \citep{golub2013matrix}: $N^3/3$ flops.  This is a one-time cost amortized
        over all $M$ paths.
  \item For each path: draw $z \sim \mathcal{N}(0, I_N)$ and compute $\nu Lz$, a
        lower-triangular matrix-vector multiply ($N^2$ flops), which gives the fBM
        path directly; then simulate the stock price path and accumulate the payoff.
\end{enumerate}

The method is exact.  The cost is $O(N^2)$ memory for $L$ (7.6~MB at $N=1000$,
which still fits in the M2's 16~MB L2 cache) and $O(MN^2)$ total work, which
dominates at practical $M$.  This method serves as the correctness baseline:
prices from the other two methods are validated against it.  For the comparison
to be fair, the baseline must do only the work it needs.  An earlier version of my
implementation multiplied by a dense copy of $L$, zeros included, at $2N^2$ flops
per path; §\ref{sec:cholesky-impl} describes the fix and its effect.

\subsection{Algorithm 2: Circulant embedding + FFT (Davies-Harte)}
\label{sec:fft}

The fBM covariance is not Toeplitz, but the covariance of the fGn
increments $\delta W_k$ is (§\ref{sec:toeplitz}). A Toeplitz matrix embeds
into a circulant; circulant matrices are diagonalized by the DFT
(§\ref{sec:circulant}).

\paragraph{Algorithm}
\begin{enumerate}
  \item Compute the fGn autocovariance $\gamma(0), \ldots, \gamma(N-1)$ and form the
        symmetric circulant first row $c \in \mathbb{R}^{2N}$.
  \item Forward FFT: $\lambda = \mathrm{FFT}(c)$; the eigenvalues are all real
        (the circulant is symmetric) and positive (§\ref{sec:circulant}).  The
        code checks the sign and throws if any is negative.
        $O(N\log N)$ flops, computed once.
  \item Pre-compute $s_j = \sqrt{\lambda_j / (2N)}$ for all $j$.
  \item For each path: draw $W \in \mathbb{C}^{2N}$ whose real and imaginary parts
        are i.i.d.\ $\mathcal{N}(0,1)$, scale $W_j \leftarrow s_j W_j$, apply inverse FFT, take the real
        part of the first $N$ entries to obtain fGn increments, cumsum to fBM.
        $O(N\log N)$ flops per path.
\end{enumerate}
FFTW plans for the forward and inverse transforms are created once before the Monte
Carlo loop and reused across all $M$ paths \citep{frigo2005design}.

No clipping is needed.  By the argument of §\ref{sec:circulant}, every eigenvalue
is positive for $H \leq \nicefrac{1}{2}$, so the method is exact on the grid at any
$N$ \citep{craigmile2003simulating}.  The smallest eigenvalue is $0.24\%$ of
$\gamma(0)$ at $H = 0.10$, $N = 252$ (§\ref{sec:stability}): small, but positive by
construction.

\subsection{Algorithm 3: Global low-rank rSVD}
\label{sec:rsvd}

Away from the diagonal $t = s$, the fBM kernel is smooth, so off-diagonal blocks of
$C$ are well-approximated at low rank (§\ref{sec:lowrank}). A global rank-$k$
approximation $C \approx U\,\mathrm{diag}(S)\,U^\top$ replaces $C$ with a compact
factor $L_k = U\,\mathrm{diag}(\sqrt{S}) \in \mathbb{R}^{N\times k}$, dropping the
per-path cost from $O(N^2)$ to $O(Nk)$.

\paragraph{Algorithm}
\begin{enumerate}
  \item Build $C \in \mathbb{R}^{N\times N}$ analytically: $O(N^2)$ flops.
  \item Compute a rank-$(k+p)$ randomized SVD of $C$ with $q = 2$ rounds of
        subspace iteration, re-orthonormalizing after every product (Algorithm~4.4
        of \citet{halko2011finding}; see §\ref{sec:poweriter}): $O(N^2 k)$ flops.
        Retain the top-$k$ singular triplets $(U, S, U^\top)$.
  \item Form the approximate Cholesky factor $L_k = U\,\mathrm{diag}(\sqrt{S})
        \in \mathbb{R}^{N\times k}$: $L_k L_k^\top \approx C$.
  \item For each path: draw $z \sim \mathcal{N}(0, I_k)$, compute $\nu L_k z$
        ($O(Nk)$ flops), cumsum to approximate fBM, simulate payoff.
\end{enumerate}

\paragraph{What it sacrifices.}
Exactness.  The sampled paths have covariance $\nu^2 C_k$ rather than $\nu^2 C$,
with the error controlled by rank $k$.  Section~\ref{sec:errorrank} shows that the
relative Frobenius error falls slowly, from $8.5\%$ at $k = 2$ to $1.2\%$ at
$k = 128$, because the spectrum of the full matrix decays only algebraically
(Figure~\ref{fig:structure}(d)).

This is a global rSVD, a single rank-$k$ approximation of the entire
$N\times N$ matrix.  It must allocate rank budget to both the smooth far-off-diagonal
blocks and the singular near-diagonal region, which drives up the required $k$
at small $H$.  A true hierarchical H-matrix keeps near-diagonal blocks dense and
compresses only far-off-diagonal blocks \citep{hackbusch2015hierarchical}.

\subsection{Other algorithms for further exploration}

\paragraph{True hierarchical H-matrix}
A recursive block-tree H-matrix decomposes $C$ into $O(\log N)$ levels of blocks:
near-diagonal blocks at each level are kept dense, while well-separated blocks at the
same level are approximated at rank proportional to their distance from the diagonal.
The resulting structure supports an H-Cholesky factorization with $O(Nk\log N)$
construction cost and $O(Nk)$ per-path cost, which is better scaling than the global rSVD
at large $N$ \citep{hackbusch2015hierarchical, bebendorf2008hierarchical}.
The complexity gain is a $\log N$ factor over Algorithm~3; it requires implementing
a recursive block-tree and the associated H-matrix arithmetic (§\ref{sec:hmatrix-future}).

\paragraph{Hybrid Scheme for rough Bergomi}
The Hybrid Scheme of \citet{bennedsen2017hybrid} achieves high RMSE reduction
over naïve Euler discretization at the same $O(N\log N)$ cost.  It does so by
exploiting the separability of the Volterra kernel $K(t,s) = (t-s)^{H-1/2}$ near
$t=s$ into a singular analytic part and a smooth remainder.  However, this kernel
governs the rough Bergomi model \citep{bayer2016pricing}, not the fBM-RFSV
model studied here: log-volatility in rough Bergomi is $\int_0^t (t-s)^{H-1/2}dW_s$,
a Volterra integral, whereas in RFSV it is simply $\nu W_t^H$.  Adopting the Hybrid
Scheme would require switching model class.

\paragraph{Quasi-Monte Carlo}
Replacing pseudo-random Gaussian draws with low-discrepancy sequences
\citep{niederreiter1992random} reduces MC standard error from $O(M^{-1/2})$ to near
$O(M^{-1})$, an improvement that can be combined with any of the three algorithms.
\section{Implementation Notes}
\label{sec:implementation}

This section documents the non-obvious design decisions in the C++ codebase
that are not fully captured by the complexity table in §\ref{sec:algorithms}.

\subsection{fGn autocovariance, not fBM covariance}
\label{sec:fgn-choice}

The circulant embedding requires a stationary covariance sequence.
fBM itself is non-stationary: the variance $\mathbb{E}[(W_t^H)^2] = t^{2H}$
grows with $t$, so row $i$ of the fBM covariance matrix $C$ is different from
row $j \neq i$.  Attempting to embed $C$ into a circulant and taking the FFT
produces a mix of positive and negative eigenvalues even for $H = \nicefrac{1}{2}$;
the method collapses entirely.

The correct input is the fGn autocovariance $\gamma(k)$ of increments
(§\ref{sec:toeplitz}):
\begin{equation}
  \gamma(k)
  = \tfrac{\Delta t^{2H}}{2}\bigl((k{+}1)^{2H} + (k{-}1)^{2H} - 2k^{2H}\bigr),
\end{equation}
which depends only on lag $k$.  The pricer uses the fGn autocovariance to
build the symmetric circulant first row
$c = (\gamma(0), \ldots, \gamma(N{-}1), 0, \gamma(N{-}1), \ldots, \gamma(1))$
and then cumsums the resulting fGn increments back into fBM log-volatility.
The bug is silent: the program runs and produces plausible-looking paths, just
with the wrong covariance structure.

\subsection{FFTW plan creation and reuse}
\label{sec:fftw-plans}

FFTW separates plan creation (analyzing the transform size and
selecting an algorithm) from execution \citep{frigo2005design}.
The implementation creates two plans once before the Monte Carlo loop:

\begin{enumerate}
  \item A forward $c2c$ plan of length $M = 2N$ to compute the circulant
        eigenvalues from the first row $c$.  This plan is executed once and
        immediately destroyed.
  \item A backward $c2c$ plan (\texttt{FFTW\_BACKWARD}) of length $M = 2N$
        to synthesise fGn increments per path.  This plan is reused across
        all $M_\mathrm{paths}$ calls to \texttt{fftw\_execute}.
\end{enumerate}

I used full complex-to-complex ($c2c$) transforms rather than the real-input
variants (\texttt{r2c}/\texttt{c2r}).  The Hermitian-symmetry bookkeeping
required for the latter, combined with the phase adjustments in the synthesis
step, introduces enough notation to obscure the algorithm; $c2c$ keeps the
implementation transparent at negligible memory cost ($2N$ vs $N{+}1$ complex
words for the eigenvalue array).

Both plans use \texttt{FFTW\_ESTIMATE}, which skips empirical tuning.  A
production deployment could substitute \texttt{FFTW\_MEASURE} or persist a
wisdom file (§\ref{sec:opt-opportunities}) for repeated calls at the same $N$.

\subsection{\texttt{std::complex<double>} and the FFTW ABI}
\label{sec:complex-cast}

A note for the interested reader / C++ developer: FFTW's C interface expects \texttt{fftw\_complex*}, which is
\texttt{double[2]}.  I allocate \texttt{std::vector<std::complex<double>>}
and pass a \texttt{reinterpret\_cast<fftw\_complex*>} of its data pointer.
This cast is valid, as the C++11 standard mandates that \texttt{complex<double>}
has the same memory layout as \texttt{double[2]}, with the real part first,
making the two representations ABI-compatible.  Using \texttt{std::vector}
rather than a raw FFTW-allocated array lets RAII manage the buffer lifetime
without a manual \texttt{fftw\_free}.

\subsection{FFT normalization convention}
\label{sec:normalization}

FFTW's backward transform is unnormalized: the output of
\texttt{fftw\_execute} on a length-$M$ array equals $M$ times the true
inverse DFT.  To obtain fGn increments with the correct covariance structure,
the per-component scaling factor is
\begin{equation}
  s_j = \sqrt{\lambda_j / M}, \qquad M = 2N,
\end{equation}
not $\sqrt{\lambda_j}$.  Applying the unnormalized backward transform to the
scaled vector $w_j = s_j(a_j + ib_j)$ (with $a_j, b_j \overset{iid}{\sim}
\mathcal{N}(0,1)$) then yields $\mathrm{Re}(\mathrm{IFFT}(w))_{0:N}$ with
covariance $\gamma(|i-j|)$, as required.

\subsection{Randomized SVD: subspace iteration}
\label{sec:poweriter}

The rSVD implementation (\texttt{rsvd.hpp}) draws a Gaussian test matrix
$\Omega \in \mathbb{R}^{N \times (k+p)}$ with oversampling $p = 5$ and
orthonormalizes the sketch, $Q = \mathrm{orth}(C\Omega)$.  It then applies $q = 2$
rounds of subspace iteration,
\begin{equation}
  Q \leftarrow \mathrm{orth}\bigl(C\,\mathrm{orth}(C^\top Q)\bigr),
  \qquad q = 2 \text{ times},
\end{equation}
where $\mathrm{orth}(\cdot)$ is a thin Householder QR.  This is Algorithm~4.4 of
\citet{halko2011finding}.  The QR after every product is what separates it from
the plain power iteration of their Algorithm~4.3.

Re-orthonormalizing is necessary here, not cosmetic.  Without it, the columns of
$(CC^\top)^q C\Omega$ are scaled by singular values raised to the power $2q+1 = 5$.
At $N = 500$, $k = 128$ the ratio $(\sigma_1/\sigma_{k+p})^5$ is
$\approx 2 \times 10^{15}$, at the limit of double precision, so the directions the
sketch is meant to capture are at risk of being lost to roundoff.  An earlier
version of my code used Algorithm~4.3.  Switching changed the Frobenius errors by
less than $0.02\%$ (relative) at the ranks tested, but it removes a failure mode at
larger $k$ or $N$.  With subspace iteration, the Frobenius error at every rank is
within about $2\%$ (relative) of the optimal Eckart--Young truncation computed from
the exact eigendecomposition.  The slow error decay in §\ref{sec:errorrank} is
therefore a property of $C$, not of the randomized algorithm.

\subsection{In-place Cholesky and the triangular mat-vec}
\label{sec:cholesky-impl}

The per-path product $\nu Lz$ needs to touch only the $N(N+1)/2$ non-zero entries
of $L$.  My first implementation copied $L$ out of \texttt{Eigen::LLT} into a dense
\texttt{MatrixXd} and computed \texttt{L * z}.  That is a general mat-vec which also
multiplies the zero upper triangle: $2N^2$ flops instead of $N^2$, and twice the
bytes read per path.  It also kept three $N \times N$ matrices alive during the
Monte Carlo loop: $C$, the factorization's internal copy, and $L$.

The current implementation factors in place, running \texttt{Eigen::LLT} on an
\texttt{Eigen::Ref} to $C$ so that $L$ overwrites the lower triangle of $C$, and
multiplies through a lower-triangular view of that storage
(\texttt{triangularView<Lower>}).  In
isolation at $N = 1000$, the triangular mat-vec takes $\approx 62\,\mu$s against
$\approx 115\,\mu$s for the dense one.  End to end, the Cholesky run at $N = 1000$,
$M = 10{,}000$ drops from $1.61$~s to $1.13$~s, and resident memory drops from
$3 \times 7.6$~MB to $7.6$~MB.  Prices are unchanged to six significant figures,
since the random draws and the arithmetic on the non-zero entries are the same.
This correction matters for the comparison in §\ref{sec:timing}: against the
dense baseline the FFT method looked $1.5\times$ faster at $N = 1000$, but against
the fair one it is only $\approx 9\%$ faster.

\subsection{Piecewise-constant volatility convention}
\label{sec:volconv}

All three pricers share a single GBM step kernel (\texttt{asian\_payoff.hpp}):
\begin{equation}
  S_i = S_{i-1}\exp\!\bigl((r - \tfrac{1}{2}\sigma_i^2)\,\Delta t
        + \sigma_i\sqrt{\Delta t}\,Z_i\bigr),
  \qquad \sigma_i = e^{\log\sigma[i]},
\end{equation}
where $\log\sigma[i] = \nu W_{t_i}^H$ is the fBM log-vol evaluated at the
end of the $i$-th interval (right-endpoint convention), and
$Z_i \overset{iid}{\sim} \mathcal{N}(0,1)$ is an independent Brownian
increment for the price process. This departs from the standard
Euler-Maruyama left-endpoint rule, which would use $\sigma_{i-1}$ (the
volatility at the start of the interval).  The departure does not constitute
a look-ahead bias in the causal sense: $W^H$ is pre-generated independently
of the price process $S$, so using $\sigma_i$ does not introduce any
dependence on future values of $S$.  It does introduce an $O(\Delta t)$
discretization error relative to the left-endpoint rule, which at $N = 252$
is negligible.  Using the same convention across all three pricers ensures
that timing and price comparisons are not confounded by volatility
discretization differences.

\subsection{Header-only pricers and ODR}
\label{sec:odr}

Each algorithm lives entirely in a \texttt{.hpp}: \texttt{cholesky.hpp},
\texttt{fft.hpp}, and \texttt{lowrank.hpp}.  Each \texttt{price()} function
is marked \texttt{inline}:
\begin{verbatim}
namespace cholesky { inline double price(...) { ... } }
\end{verbatim}
\texttt{inline} permits the same function to be defined in multiple
translation units; the linker deduplicates the definitions and keeps one
copy in the final binary.  This allows \texttt{benchmark.cpp} to include all
three headers in one translation unit without ODR violations, while each
\texttt{\_pricer.cpp} standalone executable also includes its header and
calls \texttt{price()} directly.

\subsection{Optimization opportunities not implemented}
\label{sec:opt-opportunities}

Three straightforward improvements would accelerate the MC loop without
algorithmic changes:

\paragraph{Two fGn paths per inverse FFT}
At $N = 1000$ the FFT pricer's per-path cost is dominated not by the transform but
by random number generation.  In isolation, its $4N$ Gaussian draws (real and
imaginary parts of $2N$ complex variates) take $\approx 60$--$70\,\mu$s, against
$\approx 8\,\mu$s for the size-$2N$ inverse FFT.  The real and imaginary parts of
the transformed vector each have exactly the fGn covariance, and their
cross-covariance vanishes because $\lambda_j = \lambda_{2N-j}$; I confirmed both
numerically to within Monte Carlo tolerance.  The pricer nevertheless discards the
imaginary part.  Using both would halve the Gaussian draws per fBM path, which is
the largest constant-factor saving available to Algorithm~2.

\paragraph{OpenMP path-level parallelism}
The outer \texttt{for(m)} loop over paths is embarrassingly parallel: each
path reads only the shared, read-only factor $L$ (or $L_k$, or the FFTW
eigenvalues) and accumulates independently.  A \texttt{\#pragma omp parallel
for reduction(+:payoff\_sum)} decorator with thread-local RNG state would
give near-linear multi-core scaling with no algorithmic changes.

\paragraph{FFTW wisdom files}
The current code uses \texttt{FFTW\_ESTIMATE}, which accepts the default
transform plan.  For repeated calls at a fixed $N$ (e.g., in a calibration loop
that prices the same contract many times), \texttt{FFTW\_MEASURE} or a
pre-saved wisdom file would eliminate the one-time planning overhead by
caching the optimal decomposition into FFT subproblems
\citep{frigo2005design}.

\section{Algorithm Benchmarks}
\label{sec:benchmarks}

All C++ benchmarks were run single-threaded on an Apple M2 MacBook Air (8-core, 24~GB unified
memory, rated DRAM bandwidth 100~GB/s; its largest cache is the 16~MB L2 shared by
the performance cores, and it has no L3), compiled with \texttt{-O3} and
\texttt{-march=native}, with $M=10{,}000$ paths and seven grid sizes $N$ from 64 to
4000 (64, 128, 252, 500, 1000, 2000, 4000).  Each timing is the median of five
repeats with the same seed, and the benchmark records their interquartile range.  The
range is under 2\% for most configurations (median 0.6\%) and reaches 17\% for single
points.  The machine is fanless, and a disturbed run is easy to spot this way: an
earlier run whose Cholesky timing at $N = 4000$ was 72\% too slow had an interquartile
range of 34\%, and I discarded it.  Benchmarks must also run on mains power: macOS Low
Power Mode slowed every timing by $\approx 1.8\times$ in one attempt.  Reference prices
use $M=500{,}000$ paths (see §\ref{sec:errorrank}).

\subsection{Runtime scaling}
\label{sec:timing}

Measurement setup: controlled: $M = 10{,}000$ paths; independent: $N$ from 64 to
4000, and method; dependent: total wall-clock time.

\begin{figure}[t]
  \centering
  \includegraphics[width=\linewidth]{time_vs_N.png}
  \caption{Wall-clock time vs.\ $N$ on log-log axes (left panel) with fitted power
    laws $t = c \cdot N^\alpha$, and fitted exponents $\alpha$ with prefactors $c$
    (right panel; dotted line = theoretical exponent).
    Controlled: $M=10{,}000$ paths.
    Independent: $N \in \{64, \ldots, 4000\}$, method.
    Dependent: total wall-clock time (median of 5 repeats; bars: interquartile range).}
  \label{fig:timing}
\end{figure}

Figure~\ref{fig:timing} shows the results.  A single power law describes the FFT
and rSVD methods well, with $\alpha = 1.00$ and $1.02$ ($R^2 = 1.000$ and $0.997$), but
not Cholesky: its global fit gives $\alpha = 1.48$ with $R^2 = 0.983$, because the
curve bends.  Local exponents between successive grid sizes
(Table~\ref{tab:exponents}) show what a single exponent hides.

\begin{table}[t]
  \caption{Local scaling exponents $\log(t_2/t_1)/\log(N_2/N_1)$ between successive
    grid sizes ($M = 10{,}000$).  Theoretical exponents assume the $O(MN^\alpha)$
    sampling cost dominates.}
  \label{tab:exponents}
  \centering
  \small
  \setlength{\tabcolsep}{4pt}
  \begin{tabular}{lrrrrrrr}
    \toprule
    $N$ range & 64--128 & --252 & --500 & --1000 & --2000 & --4000 & Theory \\
    \midrule
    Cholesky  & 1.05 & 1.19 & 1.28 & 1.52 & 1.75 & 2.26 & 2.0 \\
    FFT       & 1.09 & 0.87 & 1.02 & 1.02 & 1.00 & 1.04 & 1.0 \\
    rSVD ($k=32$) & 0.92 & 0.84 & 1.03 & 1.06 & 1.08 & 1.20 & 1.0 \\
    \bottomrule
  \end{tabular}
\end{table}

The Cholesky exponent climbs from about 1 to slightly above 2.  At small $N$ the
per-path cost is not the $O(N^2)$ triangular multiply but the $O(N)$ work around it:
drawing $N$ Gaussian variates, computing $2N$ exponentials for the stock price path,
and accumulating the payoff.  Above $N \approx 1000$ the mat-vec dominates, and
between $N = 2000$ and 4000 the exponent exceeds the theoretical 2, reaching 2.26,
because $L$ no longer fits in cache (§\ref{sec:memory}).  The rSVD exponent creeps
above 1 at large $N$ as its $O(N^2 k)$ construction becomes visible: 24\% of the
runtime at $N = 4000$.

The FFT method's $\log N$ factor stays invisible even over this $64\times$ range in $N$.
Figure~\ref{fig:perpath} normalizes the MC time per path by $N$.  An $O(N)$ method is
flat in this plot, an $O(N \log N)$ method rises linearly in $\log N$, and an
$O(N^2)$ method rises with slope 1.  The FFT curve is flat at 70--76~ns per time step
from $N = 64$ to $4000$.  The reason is that the transform is a small part of each
path.  In isolation at $N = 1000$, the inverse FFT takes $\approx 8\,\mu$s per transform,
while the random numbers and the price-path arithmetic, both $O(N)$, take the rest of
the $\approx 70\,\mu$s per path.

\begin{figure}[t]
  \centering
  \includegraphics[width=0.75\linewidth]{per_path_cost.png}
  \caption{Monte Carlo time per path per time step vs.\ $N$.  Flat means $O(N)$;
    rising linearly on the log-$x$ axis means $O(N\log N)$; slope 1 means $O(N^2)$.
    Controlled: $M=10{,}000$ paths.
    Independent: $N$, method.
    Dependent: MC loop time $/(MN)$.}
  \label{fig:perpath}
\end{figure}

Absolute times matter as much as exponents.  The FFT method overtakes Cholesky at
$N \approx 250$ and is $1.7\times$ faster at $N = 1000$ (0.71~s against 1.22~s) and
$6.7\times$ faster at $N = 4000$ (2.93~s against 19.66~s).  The rSVD method is fastest at
every $N$, $9.2\times$ faster than Cholesky at $N = 4000$, but it is approximate
(§\ref{sec:errorrank}).  Two implementation choices shaped these numbers
(§\ref{sec:cholesky-impl}, §\ref{sec:two-paths}): with a dense-$L$ Cholesky baseline
the FFT advantage at $N = 1000$ looked like $1.5\times$, and with a triangular baseline
but one path per transform it shrank to $9\%$.  A speedup is only as meaningful as its
baseline.

\begin{figure}[t]
  \centering
  \includegraphics[width=\linewidth]{construction_breakdown.png}
  \caption{One-time setup vs.\ per-path Monte Carlo cost breakdown.
    Light portion = construction (factorization or eigenvalue computation);
    dark portion = MC loop.  Percentages label the construction fraction.
    The FFT setup is negligible ($<1\%$); rSVD setup grows with $N$ because
    forming $L_k$ requires a full $O(N^2 k)$ sketch; the Cholesky factorization
    stays at or below 4\% of its runtime because the $O(MN^2)$ MC cost grows just as fast
    at $M = 10{,}000$.
    Controlled: $M=10{,}000$ paths.
    Independent: $N$, method.
    Dependent: wall-clock time split between construction and MC loop.}
  \label{fig:breakdown}
\end{figure}

\subsection{Memory and cache pressure}
\label{sec:memory}

Controlled: $M = 10{,}000$ paths; Independent: $N$, method; Dependent: size of the
dominant arrays resident during the MC loop and the measured lifetime peak (C++),
peak heap allocation (Python via \texttt{tracemalloc}).  The measured peak comes from
running each method once in a forked child process and reading its peak resident set
size via \texttt{wait4}, minus that of an idle child.

\begin{figure}[t]
  \centering
  \includegraphics[width=\linewidth]{memory_vs_N.png}
  \caption{Memory vs.\ $N$ for all C++ methods (left: lines = dominant array,
    hollow markers = measured lifetime peak), and the effective rate at which the
    Cholesky loop streams $L$ (right).  The dashed horizontal line marks the M2's
    16~MB L2 cache; Cholesky and the rSVD variant that retains $C$ cross it at
    $N \approx 1450$.
    Controlled: $M=10{,}000$ paths.
    Independent: $N$, method.
    Dependent: memory (MB) and effective streaming rate (GB/s), computed as bytes of
    $L$ read divided by wall time.}
  \label{fig:memory}
\end{figure}

Dense methods pay $O(N^2)$ memory.  The Cholesky factor occupies $8N^2$ bytes, 7.6~MB at
$N = 1000$ and 122~MB at $N = 4000$; the measured lifetime peaks are 11.1~MB and
142~MB, the excess being allocator and matrix-multiply workspace.  The FFT method
holds only $2N$ scale factors and two length-$2N$ complex buffers, 80~KB at
$N = 1000$, and its measured peak at $N = 4000$ is 1.2~MB.  The rSVD variant that frees
$C$ before the MC loop keeps only $L_k$ resident while sampling (0.24~MB at $N = 1000$,
$k = 32$).  Its lifetime peak, however, matches the variant that keeps $C$ (both 11.3~MB at
$N = 1000$), because $C$ must exist while $L_k$ is built.

The Cholesky loop's effective streaming rate (Figure~\ref{fig:memory}, right) shows
the memory hierarchy at work.  Dividing the bytes of the lower triangle of $L$, read
once per path, by wall time gives 14, 24, 33 and 39~GB/s at $N = 252$, 500, 1000 and
2000, then 33~GB/s at $N = 4000$.  The rate rises while the mat-vec takes over from the
$O(N)$ work, then falls once every path must stream $L$ from DRAM.  That turnover is the
signature of a memory-bandwidth-bound loop \citep{drepper2007every}, and it is why the
local exponent exceeds 2 at large $N$.  It comes later than the L2 capacity alone
predicts ($N \approx 1450$): at $N = 2000$, $L$ is already twice the L2 and the rate is
still rising, and this point varies between runs (31~GB/s in an earlier run).  At
$N = 4000$, $L$ is eight times the L2.  The figures remain effective rates, not direct
measurements of DRAM traffic; the thread-scaling test of §\ref{sec:perf} gives cleaner
evidence that the large-$N$ loop is bandwidth-bound.

The Python RFSV engine (used in Sections~\ref{sec:validation}
and~\ref{sec:sensitivity}) is vectorized across paths and so has a different
profile.  \texttt{tracemalloc} on the log-volatility generator measures a peak of
$104$ bytes per (path, time step), i.e.\ 13 float64 words per element of the
$(M, N)$ output, at every $M \geq 1000$.  That count is consistent with the working
arrays the generator holds at once, which are $2N$ wide: two real Gaussian arrays
and two complex arrays of shape $(M, 2N)$ already account for 12 words per output
element.  At $M = 10{,}000$, $N = 252$ the peak is $\approx 260$~MB, far beyond any
cache.  The measurement covers path generation only; the price-path arrays come on
top.

\begin{figure}[t]
  \centering
  \includegraphics[width=\linewidth]{memory_profile.png}
  \caption{Python RFSV engine peak heap allocation measured by \texttt{tracemalloc}.
    \textbf{(a)} Peak bytes vs.\ $N$ (log-log) for four path counts $M$; the dashed
    red line marks the M2's 16~MB L2 cache.  Scaling is $O(MN)$ (parallel to
    the $O(N)$ reference), confirming the dominant allocation grows linearly with
    both $M$ and $N$.
    \textbf{(b)} Peak bytes vs.\ total elements $M \times N$; the fitted slope
    $\approx 104$ bytes/element is 13 float64 words per output element, consistent
    with the $(M, 2N)$ real and complex working arrays of the path generator.
    Controlled: $H=0.10$, $\nu=0.52$.
    Independent: $N \in \{64, 128, 252, 512\}$, $M \in \{100, 1000, 5000, 10000\}$.
    Dependent: peak heap allocation (bytes).}
  \label{fig:memprofile}
\end{figure}

\subsection{MC convergence}
\label{sec:convergence}

Controlled: $H=0.10$, $\nu=0.52$, $\sigma_0 = 0.2$, $K=100$, $T=1$, $S_0=100$, $r=0$,
$N=252$; Independent: $M \in \{100, 250, 500, 1000, 2500, 5000,
10000, 25000\}$ and estimator (plain, control variate); Dependent: mean price and
cross-seed standard deviation over 20 independent random seeds.

\begin{figure}[t]
  \centering
  \includegraphics[width=\linewidth]{convergence.png}
  \caption{MC convergence of plain Monte Carlo and the control variate
    (§\ref{sec:cv}) on the same paths: price $\pm 1\sigma$ vs.\ number of paths $M$
    (left), and log-log plot of cross-seed standard deviation vs.\ $M$ with fitted
    slopes (right).  The dashed line is a $10^6$-path control-variate reference on the
    same $N = 252$ grid; the shaded bands are $\pm 2\hat{\sigma}/\sqrt{M}$ around it,
    with $\hat{\sigma}$ the per-path standard deviation of each estimator in the
    reference run.
    Controlled: $N=252$, $H=0.10$, $\nu=0.52$, $\sigma_0 = 0.2$.
    Independent: $M$ (left and right); 20 seeds (right).
    Dependent: mean price $\pm 1\sigma$ and log-log slope.}
  \label{fig:convergence}
\end{figure}

Figure~\ref{fig:convergence} shows the results.  The fitted log-log slope of plain
Monte Carlo is $-0.488$ ($R^2 = 0.986$), in good agreement with the theoretical $-0.5$
from the central limit theorem \citep{glasserman2004monte}.  The reference on the same
grid is $5.3217 \pm 0.0018$ from $10^6$ paths with the control variate (plain:
$5.3252 \pm 0.0095$ from the same paths), and every seed-mean lies within its band.
The payoff standard deviation is $\sigma_\mathrm{payoff} = 9.46$, so $M = 10{,}000$ paths
give a standard error of $\approx 0.095$, $1.8\%$ of the price.  To bring the 95\%
confidence half-width below $1\%$ of the price would take
$M \approx (2\sigma_\mathrm{payoff}/0.053)^2 \approx 130{,}000$ paths.

The control variate keeps the $M^{-1/2}$ rate and lowers the constant.  Its per-path
standard deviation is $1.82$, a variance reduction of $27.1$, so the same $1\%$
half-width needs only $\approx 4{,}700$ paths.  Its fitted slope, $-0.462$
($R^2 = 0.965$), is noisier than the plain one because each cross-seed standard
deviation is itself estimated from only 20 seeds; the per-seed variance reductions
range from 14 to 57 around the per-path value of 27.  Where differences between
configurations matter, I also use common random numbers, which makes those
differences far more precise than the absolute prices.

Two earlier versions of this study failed this test, and both failures are
instructive.  The first estimated $\sigma_\mathrm{payoff}$ from the spread of only five
seed-level prices, and fitted a slope of $-0.71$: a standard deviation estimated from
five samples has roughly 35\% relative error.  The second used $\sigma_0 = 1$, and its
payoff was so heavy-tailed that the sample standard deviation grew with the sample
size, from 88 over $2 \times 10^5$ paths to 118 over $10^6$, with a single path carrying a
third of the sample variance.  The Monte Carlo error bars of that parameterization
were not meaningful; at $\sigma_0 = 0.2$ the standard deviation of 50{,}000-path batches
ranges only from 9.3 to 9.6.

\subsection{Numerical stability}
\label{sec:stability}

\begin{figure}[t]
  \centering
  \includegraphics[width=\linewidth]{stability_report.png}
  \caption{Numerical stability of the three methods.
    \textbf{(a)} FFT: smallest circulant eigenvalue relative to $\gamma(0)$ vs.\ Hurst
    exponent $H$ ($N=252$ and $N=1008$); no eigenvalue is negative anywhere.
    \textbf{(b)} rSVD: condition number $\kappa(L_k)$ and Frobenius error vs.\ rank
    $k$ ($N=252$, $H=0.10$).
    \textbf{(c)} Cholesky: condition number $\kappa(C)$ vs.\ $N$ ($H=0.10$) with
    fitted power law.}
  \label{fig:stability}
\end{figure}

\paragraph{(a) FFT eigenvalue positivity: }
Figure~\ref{fig:stability}(a) sweeps $H \in [0.05, 0.501]$ at $N = 252$ and
$N = 1008$ and plots the smallest circulant eigenvalue relative to $\gamma(0)$.  No
eigenvalue is negative at any $H$ or $N$ tested, so nothing is clipped and the FFT
sampler is exact.  The minima agree with the closed form~\eqref{eq:lambda-min} to
every printed digit: $0.24\%$ at $H = 0.10$, $N = 252$ and $0.08\%$ at $N = 1000$,
always attained at $\lambda_0$.  The margin collapses as $H \to 0$, falling to
$0.07\%$ at $H = 0.05$, $N = 252$, but it is positive by construction and many orders
of magnitude above double-precision roundoff.  At $H = \nicefrac{1}{2}$ the
increments are white noise, every eigenvalue equals $\gamma(0)$, and the ratio is $1$.
Consistent with exactness, the 500{,}000-path FFT and Cholesky prices in
§\ref{sec:errorrank} agree to within MC noise.

\paragraph{(b) rSVD condition number: }
Figure~\ref{fig:stability}(b) plots $\kappa(L_k) = \sqrt{S_\mathrm{max}} /
\sqrt{S_\mathrm{min}}$ as a function of rank $k$ at $H = 0.10$, $N = 252$.  The
condition number grows from $8.2$ at $k = 8$ to $23.3$ at $k = 64$: as we include
smaller singular values (the denominator shrinks), $\kappa$ grows.  A condition number
of 23 is well within the range of float64 arithmetic ($\epsilon \approx 10^{-16}$,
so $\kappa \cdot \epsilon \approx 2 \times 10^{-15}$) and does not cause numerical
issues in path generation.

\paragraph{(c) Cholesky covariance conditioning: }
Figure~\ref{fig:stability}(c) shows the condition number
$\kappa(C) = \lambda_\mathrm{max} / \lambda_\mathrm{min}$ against $N$.  At $N = 252$, $\kappa(C) \approx 789$; at $N =
1000$, $\kappa(C) \approx 4 \times 10^3$.  The fitted power law is
$\kappa \approx 1.08\,N^{1.19}$: the covariance matrix becomes increasingly
ill-conditioned as more time steps are added, but only polynomially.  This growth is expected: adding time steps increases
correlations between adjacent points, pushing $\lambda_\mathrm{min}$ toward zero.
The matrix remains positive-definite and Cholesky-factorizable throughout the tested
range.  \texttt{Eigen::LLT} fails when floating-point roundoff destroys
positive-definiteness, which occurs when $\kappa(C) \cdot \epsilon_\mathrm{mach}
\geq 1$, i.e., $\kappa(C) \geq 10^{16}$.  At $N = 1000$ where
$\kappa(C) \approx 4 \times 10^3$, this threshold is twelve orders of magnitude
away; the fit would not reach it until $N \sim 10^{12}$.

The slow singular value decay visible in Figure~\ref{fig:structure}(d) is a
related but distinct phenomenon: a large $\kappa(C)$ means the spectrum spans
many decades, but slow decay (many non-negligible eigenvalues) is what
actually hinders low-rank approximation.  A matrix can have an enormous
condition number and still be perfectly approximated at low rank if the
eigenvalues drop off rapidly after the first few; in the fBM case, rough $H$
implies both a large condition number and a slow decay rate.  As
Figure~\ref{fig:structure}(d) shows, that slow decay belongs to the full matrix;
its well-separated off-diagonal blocks still compress to rank $\leq 3$.

\subsection{Low-rank approximation quality}
\label{sec:errorrank}

Controlled: $N = 500$, $M = 100{,}000$ paths per rank (rSVD) and $M = 500{,}000$
paths per exact method (reference); Independent: rSVD target rank
$k \in \{2, 4, 8, 16, 32, 64, 128\}$; Dependent: Frobenius error, fraction of path
variance lost, and relative price error.

\begin{figure}[t]
  \centering
  \includegraphics[width=\linewidth]{error_vs_rank.png}
  \caption{Low-rank rSVD approximation quality vs.\ rank $k$: the Frobenius error
    $\|C - C_k\|_F/\|C\|_F$, the fraction of path variance discarded
    $\mathrm{tr}(C - C_k)/\mathrm{tr}(C)$ (both exact), and the relative underpricing
    against the reference with $\pm 2$ standard-error bars.
    Controlled: $N=500$, $H=0.10$, $M = 100{,}000$ paths per rank.
    Independent: rank $k$.
    Dependent: relative error (\%).}
  \label{fig:errorrank}
\end{figure}

Figure~\ref{fig:errorrank} shows results.  The reference price
$p_\mathrm{ref} = 5.3144$ is the average of a 500{,}000-path Cholesky run ($5.3195$)
and a 500{,}000-path FFT run ($5.3093$); averaging two independent exact methods
avoids privileging either as ground truth.  Each run has a standard error of
$\approx 0.013$, and their difference of $0.010$ is about half the standard error of a
difference, consistent with both methods being exact.

With $100{,}000$ paths per rank the price error is resolved, and it is a systematic
underpricing: $-8.1\%$ at $k = 2$, $-6.8\%$ at $k = 8$, $-4.3\%$ at $k = 32$, and $-3.9\%$
at $k = 128$, each many times the standard error of $\approx 0.56\%$.  The Frobenius
error suggests a much better approximation than these prices do: it falls from
$8.5\%$ to $1.2\%$ over the same ranks.

The key pitfall is that the Frobenius norm is the wrong measure of what matters for
pricing.  It is dominated by the few large eigenvalues of $C$, which a rank-$k$
approximation captures.  The option price depends instead on how much variance the
simulated paths carry, and a rank-$k$ truncation discards the long algebraic tail of
small eigenvalues, which hold the short-scale, rough part of the path.  The fraction
of path variance lost, $\mathrm{tr}(C - C_k)/\mathrm{tr}(C)$, is $37.9\%$ at $k = 2$,
$20.8\%$ at $k = 32$, and still $13.2\%$ at $k = 128$; at the first time step a rank-128
approximation keeps only 60\% of the variance.  Less volatility means a lower Asian
price, and the price bias decays slowly in the same way as the variance lost.  For
low-rank samplers of rough processes, the trace error should be reported alongside,
or instead of, the Frobenius error.

\section{Model Validation}
\label{sec:validation}

Three experiments probe the RFSV pricer: (i) the implied volatility smile
against live SPY option chain data; (ii) a comparison with constant-volatility
baselines, and the roughness premium with and without variance matching; and (iii)
a parameter sensitivity sweep over $H$ and $\nu$.  Each uses the Python RFSV engine
(\texttt{data/rfsv\_model.py}), which implements the circulant-FFT path
generator with the same normalization as the C++ FFT pricer.  Every comparison
between configurations (different $H$, $\nu$ or $K$) uses common random numbers, so
differences are measured far more precisely than the absolute prices.

\subsection{Implied volatility smile}
\label{sec:iv}

\begin{figure}[t]
  \centering
  \includegraphics[width=\linewidth]{validate_iv.png}
  \caption{RFSV vs market implied volatility smile for SPY, for the two nearest
    weekly expirations in a snapshot taken on October 2, 2026 (spot \$771.09).
    Scatter: market IV computed by Black--Scholes inversion of mid-quote
    (\texttt{brentq}, $r = 0$, filtered to strikes in $[0.60S, 1.40S]$ with
    valid bid/ask spread); the $y$-axis is capped at $40\%$, and the 39 and 42
    deep-ITM market IVs above it are marked as off scale.
    Line: RFSV IV, same inversion applied to Monte Carlo European call prices,
    with common random numbers across strikes and the sample mean of $S_T$
    matched to the forward.
    The RFSV log-vol level $\mu_0 = \log\sigma_\mathrm{eff}$ is calibrated per expiration to
    match the ATM price (not IV): Brent's method on the price residual
    $p_\mathrm{RFSV} - p_\mathrm{mkt}(\sigma_\mathrm{ATM})$.
    The shaded region marks deep-ITM strikes where no simulated path finishes
    below $K$, so the Monte Carlo time value is zero and the IV is undefined.
    Controlled: $H = 0.10$, $\nu = 0.52$, $r = 0$, $M = 20{,}000$,
    $N = 63$ steps per year (at least 5 steps).
    Independent: strike $K$, expiration $T$.
    Dependent: implied volatility $\sigma_\mathrm{IV}$.}
  \label{fig:iv}
\end{figure}

Figure~\ref{fig:iv} shows the results.  For this comparison the base level is not
fixed at $\sigma_0 = 0.2$ but calibrated per expiration, $\mu_0 = \log\sigma_\mathrm{eff}$,
so that the RFSV ATM price matches the market ATM price.  This aligns the curves at
the money ($12.1\%$ and $13.1\%$ for the two expirations, with
$\sigma_\mathrm{eff} = 0.112$ and $0.119$), so the comparison concerns smile shape
(slope and curvature), not the vol level.

The model smile has clear curvature but no skew.  For the nearer expiration
($T = 0.019$) it has its minimum at the money and rises to $\approx 17.5\%$ at
$K/S = 0.91$ and $\approx 15\%$ at $K/S = 1.06$.  The SPY smile is strongly skewed
instead: above $40\%$ at $K/S = 0.91$, falling to a minimum of $\approx 10.5\%$ near
$K/S = 1.02$, and recovering to $\approx 16\%$ at $1.06$.  The later expiration
($T = 0.038$) has the same shape.  In the out-of-the-money call wing the two curves
roughly meet, so the vol-of-vol $\nu = 0.52$ produces about the right amount of
curvature there.  Below the money the model misses the skew entirely.  This is
expected: the engine draws the price shocks independently of the volatility driver,
so the spot-vol correlation is $\rho = 0$, and with $\rho = 0$ a stochastic volatility
model produces a smile that is symmetric in log-moneyness.  The steep short-dated
skew that rough volatility is known for \citep{gatheral2018volatility} requires
$\rho < 0$ (§\ref{sec:rho-future}).

\paragraph{A Monte Carlo pitfall: forward error masquerading as skew.}
An earlier run of this experiment (an October 1 snapshot with $\nu = 0.30$) used
$M = 3{,}000$ paths with common random numbers but no correction, and its RFSV curve rose steadily with strike, from
$13.0\%$ to $15.7\%$ across $0.98 \leq K/S \leq 1.05$.  That curve looks like an
upward skew, which a $\rho = 0$ model cannot produce.  The cause was sampling
error in the forward.  The sample mean of $S_T$ was $0.048\%$ below $S_0$, and
because every strike shares the same paths, that error acts like a shifted
forward and tilts the entire IV curve; it also biases the calibrated $\mu_0$.  A
400{,}000-path run of the same model gives an approximately symmetric smile with its minimum at
the money ($14.57\%$ at $K/S = 1.00$, against $14.79\%$ at $0.98$ and $14.70\%$
at $1.02$).  The current script rescales the simulated $S_T$ so that its sample
mean equals the forward exactly, a standard moment-matching correction
\citep{glasserman2004monte}, and uses $M = 20{,}000$ paths.  At $N \approx 5$
time steps per path the whole experiment runs in a few seconds.

\paragraph{Limitations}
I set $r = 0$ and ignore SPY's dividend yield.  For near-the-money options this
shifts IV levels slightly and leaves the shape comparison largely unaffected.
For deep in-the-money calls the effect is large: the carry in an ITM call price
is misread as time value, which inflates those market IVs (many above $40\%$),
so I do not interpret the low-strike wing.  SPY options are American-style; for
calls on a low-dividend ETF the early-exercise premium is small, and I treat
them as European.  The comparison is a single-day snapshot; market IV smiles
vary substantially day-to-day.  Brent's method stops at a tolerance of $0.01$ in
$\mu_0$, about $1\%$ in $\sigma_\mathrm{eff}$.

\subsection{Lévy benchmark and roughness premium}
\label{sec:levy}

\begin{figure}[t]
  \centering
  \includegraphics[width=\linewidth]{validate_asian.png}
  \caption{Asian option prices vs strike.
    \textbf{(a)} Price vs $K$: Lévy (1992) approximation and exact GBM by Monte
    Carlo at $\sigma = 0.2$ (constant vol), and RFSV at four $H$ values
    ($\nu = 0.52$, $\sigma_0 = 0.2$); each curve is the mean of 10 seeds of 10{,}000
    paths, with the same seeds for every curve.
    \textbf{(b)} Roughness premium $\hat{p}(H{=}0.10) - \hat{p}(H{=}0.50)$ vs $K$, at
    fixed $\nu = 0.52$ (red) and with $\nu$ at $H = 0.5$ chosen to match the expected
    integrated variance (blue); error bars are $\pm 1$ paired standard error across
    seeds.
    Controlled: $S_0 = 100$, $T = 1$, $\sigma_0 = 0.2$, $r = 0$, $N = 252$.
    Independent: $K$ (both panels); $H$ and $\nu$.
    Dependent: arithmetic Asian call price.}
  \label{fig:levy}
\end{figure}

Figure~\ref{fig:levy} shows the results.

\paragraph{Constant-volatility baselines: }
When $\nu \to 0$, $\sigma_t \to \sigma_0$, so the RFSV dynamics reduce to
constant-volatility GBM with $\sigma = 0.2$, which the Lévy approximation is designed
for \citep{levy1992pricing}.  At this volatility it is accurate: at the money it gives
$4.625$, against $4.606 \pm 0.027$ for exact discrete GBM by Monte Carlo
(100{,}000 paths).  All four RFSV curves lie above both baselines, ordered by $H$, as
stochastic volatility and Jensen's inequality predict.  At the money the RFSV prices
are $5.445$, $5.302$, $4.979$, and $4.834$ for $H = 0.05$, $0.10$, $0.30$, and $0.50$, so
rough stochastic volatility adds $0.70$ to the GBM price at $H = 0.10$.  (An earlier
version of this experiment used $\sigma = 1$, where the Lévy approximation overprices
exact GBM by about 6\%, enough to put the Lévy curve above the RFSV curves; it is
reliable only at realistic volatility.)

\paragraph{Roughness premium: }
At fixed $\nu = 0.52$, the premium $\hat{p}(H{=}0.10) - \hat{p}(H{=}0.50)$ is positive at
every strike, rising from $0.178$ at $K = 80$ to $0.468 \pm 0.008$ at the money and
falling to $0.319$ at $K = 120$.  But this compares more than roughness.  At fixed
$\nu$, $\text{Var}(\log\sigma_t) = \nu^2 t^{2H}$ is larger for small $H$ when $t < 1$, and
the expected integrated variance
\begin{equation}
  \label{eq:intvar}
  \mathbb{E}\!\left[\int_0^T \sigma_t^2\,dt\right]
  = \sigma_0^2 \int_0^T e^{2\nu^2 t^{2H}}\,dt
\end{equation}
is $18.6\%$ higher at $H = 0.10$ than at $H = 0.50$ ($1.574\,\sigma_0^2$ against
$1.327\,\sigma_0^2$).  To isolate roughness, I solve~\eqref{eq:intvar} for the $\nu$ at
$H = 0.5$ that gives the same integrated variance as $(H, \nu) = (0.10, 0.52)$, which
is $\nu = 0.651$, and price again with the same random numbers.  The variance-matched
premium is $0.330 \pm 0.009$ at the money, so about 70\% of the premium survives: most
of the price increase comes from path roughness itself, not from higher total
variance.  Both premia peak at the money, where the call payoff is most sensitive to
the shape of the volatility path.

\paragraph{Limitations: }
Matching expected integrated variance is one reasonable normalization, not the only
one; matching $\int_0^T \text{Var}(\log\sigma_t)\,dt$ instead gives a slightly different
$\nu$ (§\ref{sec:variancenorm}).  The standard errors are paired across 10 seeds, so
they reflect how consistently the premium appears under common random numbers; the
absolute prices carry the larger standard error of $\approx 0.03$.

\subsection{Parameter sensitivity: }
\label{sec:sensitivity}

\begin{figure}[t]
  \centering
  \includegraphics[width=\linewidth]{sensitivity_surface.png}
  \caption{ATM Asian call price vs model parameters.
    \textbf{(a)} Heatmap: price vs $(H, \nu)$ at $K = S_0 = 100$ (viridis
    colormap; low H = rough = lower on the $y$-axis because the axis is
    inverted so H increases upward).
    \textbf{(b)} Price vs $H$ for each $\nu$; curves decrease left-to-right
    and separate vertically with $\nu$.
    All cells share the same random numbers, so differences between cells are
    precise even though each absolute price has an MC standard error of
    $\approx 0.1$.
    Controlled: $S_0 = 100$, $K = 100$, $T = 1$, $\sigma_0 = 0.2$, $r = 0$, $N = 252$,
    $M = 10{,}000$.
    Independent: $H$, $\nu$.
    Dependent: ATM Asian call price.}
  \label{fig:sensitivity}
\end{figure}

\begin{figure}[t]
  \centering
  \includegraphics[width=\linewidth]{sensitivity_strike.png}
  \caption{Asian call price vs strike $K$ for varying $H$ at $\nu = 0.52$, with
    common random numbers across all curves.  Curves decrease monotonically with
    strike and are ordered by $H$ at every strike; the spread between $H = 0.05$
    and $H = 0.50$ peaks at the money.
    Controlled: $\nu = 0.52$, $\sigma_0 = 0.2$, $S_0 = 100$, $T = 1$, $r = 0$,
    $N = 252$, $M = 10{,}000$.
    Independent: $K$, $H$.
    Dependent: Asian call price.}
  \label{fig:sensitivitystrike}
\end{figure}

Figures~\ref{fig:sensitivity} and~\ref{fig:sensitivitystrike} show the
results:

\begin{itemize}
  \item \textbf{Price decreases with $H$, and more so at large $\nu$:} at
        $\nu = 0.52$ the ATM price falls from $5.373$ at $H = 0.05$ to $4.735$ at
        $H = 0.50$.  The gap between $H = 0.10$ and $H = 0.50$ is only $0.020$ at
        $\nu = 0.10$ but $0.937$ at $\nu = 0.70$: without vol-of-vol there is no
        volatility path whose roughness could matter.  As §\ref{sec:levy} shows,
        most of this effect is roughness rather than added variance.

  \item \textbf{Price increases with $\nu$:} at $H = 0.10$ the ATM price rises from
        $4.540$ at $\nu = 0.10$ to $5.869$ at $\nu = 0.70$.  Larger vol-of-vol
        amplifies $\sigma_t$ variation, and the increase is convex in $\nu$ because
        $\sigma_t = \sigma_0 e^{\nu W_t^H}$ enters exponentially.

  \item \textbf{The $H$ effect peaks at the money:} the spread between $H = 0.05$
        and $H = 0.50$ is $0.27$ at $K = 80$, $0.64$ at $K = 100$, and $0.41$ at
        $K = 120$ (Figure~\ref{fig:sensitivitystrike}), matching the strike profile
        of the premium in Figure~\ref{fig:levy}(b).
\end{itemize}

At the model parameters $H = 0.10$, $\nu = 0.52$ the ATM price in the grid is $5.223$,
one standard error ($\approx 0.095$) below the $10^6$-path reference of $5.325$
(§\ref{sec:convergence}); the grid's absolute level carries that single run's MC error,
while its differences do not.

\paragraph{Limitations: }
All grid points share one random seed.  This makes the trends across the grid
reliable, but the whole surface carries a common MC offset with a standard error of
$\approx 0.1$ price units.  The sensitivity surfaces are computed on a $6 \times 4$ grid
($H \times \nu$) at a fixed $N = 252$; finer grids and higher $M$ are feasible but
were not run for this report.

\section{Future work}
\label{sec:futurework}

The three algorithms implemented here exploit three structural properties of
the fBM covariance matrix: positive-definiteness (Cholesky), stationarity of
increments (circulant FFT), and smooth off-diagonal low-rank structure (global
rSVD).

\subsection{True hierarchical H-matrix}
\label{sec:hmatrix-future}

The global rSVD (Algorithm~3) approximates the entire $N \times N$
matrix $C$ at rank $k$.  This forces the rank budget to cover both the
smooth far-off-diagonal region and the singular near-diagonal region, where
$C(t,s) \sim |t-s|^{2H}$ with $H = 0.10$ is Hölder-$0.2$ continuous and
barely above $C^0$ regularity.  The global approximation cannot specialize, which is
why the spectrum of the full matrix decays slowly at small $H$ even though its
off-diagonal block compresses to rank $\leq 3$ (Figure~\ref{fig:structure}(d)).

A true hierarchical H-matrix separates these concerns by a recursive
block-tree decomposition of $C$ \citep{hackbusch2015hierarchical,
bebendorf2008hierarchical}.  At each level of the tree:
\begin{itemize}
  \item Near-diagonal blocks (time intervals that overlap or touch)
        are stored dense.  The near-diagonal singularity is not low-rank, so
        no rank budget is wasted trying to compress it.
  \item Far-off-diagonal blocks (well-separated time intervals) lie
        entirely in the smooth region $t \neq s$ and are approximated at low
        rank proportional to the distance from the diagonal.
\end{itemize}
The result is an H-Cholesky factorization with $O(Nk\log N)$ construction
cost and $O(Nk\log N)$ per-path cost (H-matrix triangular solve traverses
the block tree), versus the global rSVD's $O(Nk)$ per-path cost.  The
H-matrix advantage lies in construction: $O(Nk\log N)$ vs $O(N^2 k)$ for the
global rSVD sketch, a factor of $N/\log N$ for $k \ll N$.

For the rough case $H = 0.10$, the off-diagonal block needs rank $3$ rather than
$2$ for $1\%$ accuracy (Figure~\ref{fig:structure}(d)): slightly more than at
$H = 0.50$, but far below the ranks the global approximation uses.  The H-matrix
benefit over the global rSVD grows as $N$ increases beyond the current test range.  At $N = 252$,
$O(N^2 k) = O(252^2 \cdot 64) \approx 4 \times 10^6$ flops for construction;
the H-matrix would reduce this by roughly $\log_2 252 \approx 8$ for equal
per-block rank, which becomes worthwhile at $N \geq 1000$.

\subsection{Hybrid Scheme for rough Bergomi}
\label{sec:hybrid-future}

The Hybrid Scheme of \citet{bennedsen2017hybrid} achieves more than $80\%$
RMSE reduction over naïve Euler discretization at $O(N \log N)$ cost by
exploiting the separability of the Volterra kernel $K(t,s) = (t-s)^{H-1/2}$
near $t = s$: the singular part is handled analytically and the smooth
remainder is convolved via FFT.  This is a structurally superior simulation
method for the rough Bergomi model of \citet{bayer2016pricing}, in
which log-volatility is a Volterra integral $\int_0^t (t-s)^{H-1/2}\,dW_s$.

The RFSV model studied here uses a different representation: $\log\sigma_t =
\nu W_t^H$, where $W_t^H$ is a standard fBM (not a Volterra integral with
the Bergomi kernel).  The two models are different processes with different
correlation structures, though both have Hurst exponent $H$ controlling
roughness.  Adopting the Hybrid Scheme would require switching to the rough
Bergomi model class, which has distinct option pricing properties and
calibration procedures.  The structural insight of separating the kernel into
a singular analytic part and a smooth FFT-able remainder is directly
relevant to the fBM kernel $C(t,s)$ and motivates the hierarchical treatment
in §\ref{sec:hmatrix-future}.

\subsection{Other variance normalizations}
\label{sec:variancenorm}

Section~\ref{sec:levy} isolates roughness by matching the expected integrated
variance~\eqref{eq:intvar} across $H$, which leaves about 70\% of the premium.  Other
normalizations are defensible and would shift that share.  Matching the integrated
log-volatility variance $\int_0^T \nu^2 t^{2H}\,dt$ instead gives the closed form
\begin{equation}
  \nu(H) = \nu_0\sqrt{\frac{2H+1}{2}\,T^{1-2H}},
\end{equation}
with $\nu_0$ fixed at the Brownian reference ($H = \nicefrac{1}{2}$).  From
$(H, \nu) = (0.10, 0.52)$ this gives $\nu = 0.671$ at $H = 0.5$, against $0.651$ from
matching~\eqref{eq:intvar}.  A third option holds the at-the-money implied volatility
fixed, which is closer to how a trader would compare models.  Reporting the premium
under all three would show how robust the roughness share is.

\subsection{Spot-volatility correlation}
\label{sec:rho-future}

The engine draws the price shocks $Z$ independently of the fBM driving
volatility, so the spot-vol correlation is $\rho = 0$.  As §\ref{sec:iv} shows,
this leaves the model unable to produce the downward skew of equity smiles.
Introducing $\rho \neq 0$ requires simulating $W^H$ jointly with the Brownian motion
that drives the price.  The dense Cholesky sampler can absorb this directly by
factoring the $2N \times 2N$ joint covariance, which is explicit when $W^H$ is
written as a Volterra integral of that Brownian motion; the circulant embedding
would need to be extended to the joint process.  This is the single change most
needed before the IV comparison can test the model rather than its
simplifications.

\subsection{Quasi-Monte Carlo path generation}
\label{sec:qmc-future}

All three algorithms use pseudo-random Gaussian draws, giving a standard MC
error of $O(M^{-1/2})$.  Replacing these with low-discrepancy sequences
\citep{niederreiter1992random} reduces the effective error rate to
approximately $O(M^{-1}(\log M)^N)$, close to $O(M^{-1})$ in practice
\citep{glasserman2004monte}.  This improvement is orthogonal to the
choice of covariance simulation method: any of the three algorithms can be
combined with QMC by replacing the standard normal draws.

In $N = 252$ dimensions the low-discrepancy property is strongest for the
first few coordinates.  Brownian bridge dimension ordering concentrates the
largest variance components in the lowest-index dimensions, ensuring the
low-discrepancy sequence acts on the high-information components first
\citep{joy1996quasi}.  At $M = 10{,}000$, the MC standard error of
$\approx 0.095$ (§\ref{sec:convergence}) would fall to $\approx 0.001$ under
ideal $O(M^{-1})$ QMC, a $100\times$ reduction at no additional computational
cost, improving the statistical resolution of the roughness premium and
parameter sensitivity results.


\medskip

\bibliography{bibliography}

\appendix
\section*{Reproducibility}
All source code, benchmark scripts, and plotting code are available at
\url{https://github.com/t-flora/fbm-fast-pricers}.  Results can be reproduced
end-to-end by running \texttt{./run\_pipeline.sh} from the project root.  A manifest
recording the git commit, parameters, and wall time is written to
\texttt{plots/runs/\textit{TIMESTAMP}/manifest.txt} for each run.

\section*{Acknowledgements}
Portions of this report were drafted and edited with the assistance of Claude Sonnet 4.6
(Anthropic, 2026), Claude Opus 5.5 (Anthropic, 2026), and Gemini 3.1 Pro (Google, 2026).

\end{document}
